% explainer-source-hash: 758aa0e1fc4aa54b67ece6e4982367fb6dc2c725dd0193a8e15921f8ec2b6aa1
% explainer-source-commit: a16d032e0874aa3ffd34322c529bf16d461e5936
% explainer-generated: 2026-07-16
% Regenerate with /explain-all (see WIRING.md). Do not edit this stamp by hand:
% it is how the toolchain knows whether this document still matches the code.
\documentclass[11pt,a4paper]{article}
\input{preamble}

\title{\vspace{-1.5cm}\textbf{N8N Automations}\\[0.3em]
\large The Brief\\[0.2em]
\normalsize 59 business chores, each built as a flowchart that then runs itself}
\author{}
\date{}

\begin{document}
\maketitle
\thispagestyle{empty}

\section{The project in one page}

Fifty-nine business processes, each built as a flowchart inside a tool called n8n,
exported to text files, and documented one folder at a time. Scraping leads, chasing
enquiries, sending invoices, writing posts, rendering videos, answering questions from a
company's own documents.

There is no program here to compile and no test suite to run. The unit of work is a
\emph{workflow}: a chain of steps wired together on a canvas, started by an event, that
moves data between web services and, more often than not, hands something to an AI model
along the way.

\begin{keyidea}{The three numbers that describe this repository}
\begin{center}
\begin{tabular}{@{}rp{10.4cm}@{}}
\textbf{59}      & automation folders, each with a README, a credential manifest, a licence and an export \\[2pt]
\textbf{80}      & exported workflow files (six folders hold multi-workflow systems) \\[2pt]
\textbf{1{,}645} & nodes, of which 186 are sticky notes documenting the other 1{,}459 \\
\end{tabular}
\end{center}
All three were verified by parsing the files, not by counting the README's bullets. The
headline claim of \textbf{59 automations is accurate}.
\end{keyidea}

\section{What n8n is, for a reader who has never seen it}

Take a real chore: \emph{when a deal is marked won in our customer database, create a
customer in our payment processor and invoice them}. You can do it by hand in four
minutes, or write and host a program for a day. n8n is the third option: drag a box for
each step, wire the boxes together, and it runs the flowchart for you. It already knows
how to talk to hundreds of services.

\begin{figure}[H]
\centering
\resizebox{\textwidth}{!}{
\begin{tikzpicture}[node distance=7mm]
  \node[extbox] (world) {Outside world\\\scriptsize a deal is marked "won"};
  \node[box, right=13mm of world, fill=warnred!12, draw=warnred] (trig) {\textbf{Trigger}\\\scriptsize webhook};
  \node[box, right=10mm of trig] (fetch) {Node\\\scriptsize fetch the deal};
  \node[box, right=10mm of fetch] (cust) {Node\\\scriptsize create customer};
  \node[box, right=10mm of cust] (inv) {Node\\\scriptsize send invoice};
  \draw[flow] (world) -- (trig);
  \draw[flow] (trig) -- node[lbl,above]{data} (fetch);
  \draw[flow] (fetch) -- node[lbl,above]{data} (cust);
  \draw[flow] (cust) -- node[lbl,above]{data} (inv);
  \node[store, below=11mm of fetch] (cred) {credential\\store};
  \draw[dflow] (cred) -- (fetch);
  \draw[dflow] (cred) -- (cust);
  \draw[dflow] (cred) -- (inv);
  \node[lbl, below=1mm of cred] {logins live here, never in the flowchart};
\end{tikzpicture}}
\caption{The whole vocabulary. Each box is a \textbf{node}. Each arrow is a
\textbf{connection}. The box that starts it is the \textbf{trigger}. The logins are
\textbf{credentials}, referenced from a separate store. The chain is a
\textbf{workflow}.}
\end{figure}

\begin{keyidea}{Dragging boxes removes the typing, not the thinking}
Everything hard about software survives the move to a canvas: state, failure, retries,
duplicate events, data shape mismatches, rate limits. Section~\ref{sec:weak} is a list of
places in this repository where that bill came due, and every item on it is an ordinary
software bug that would have been equally possible in handwritten code.
\end{keyidea}

\section{What the collection is made of}

\begin{figure}[H]
\centering
\begin{tikzpicture}
\begin{axis}[
  width=13cm, height=5.6cm,
  xbar, xmin=0,
  bar width=8pt, y=11pt,
  enlarge y limits=0.08,
  xlabel={\small nodes across all 80 exports},
  symbolic y coords={code,if,lmChatGoogleGemini,agent,googleSheets,httpRequest,set,stickyNote},
  ytick=data,
  nodes near coords, nodes near coords style={font=\tiny},
  tick label style={font=\scriptsize}, label style={font=\small},
  every axis plot/.append style={fill=accent!55, draw=inkblue},
  axis lines*=left, xmajorgrids, grid style={linegrey},
]
\addplot coordinates {
  (186,stickyNote) (136,set) (128,httpRequest) (124,googleSheets)
  (109,agent) (82,lmChatGoogleGemini) (80,if) (71,code)
};
\end{axis}
\end{tikzpicture}
\caption{The eight most common node types, and a fair portrait of the whole collection:
moving data between web services and an LLM, glued together with spreadsheets, and
unusually well commented.}
\end{figure}

Four readings fall straight out of that chart, and each is worth a sentence.

\textbf{Sticky notes are the most common node.} A sticky note does nothing at run time; it
is a comment on the canvas. One node in nine here exists only to explain the others. That
is a genuinely unusual signal of care.

\textbf{Google Sheets is the database} (124 nodes). It is the lead list, the job queue, the
state store, the log, and the client's review screen, all at once.

\textbf{128 hand-written API calls.} Wherever n8n had no purpose-built node, the author read
the vendor's docs and constructed the raw call. This is the single most important thing to
notice about the collection: it is where the drag-and-drop tool runs out and the
engineering starts.

\textbf{AI is structural, not decorative.} 109 agent nodes, 82 Gemini models, 37 OpenAI
nodes. About a fifth of the repository.

\section{The seven families}

Fifty-nine is too many to list, and they are not 59 unrelated things. Grouping them by
what they do (this grouping is this document's own reading; the repository does not group
its folders) gives seven families whose architecture is largely predictable from the
family alone.

\begin{figure}[H]
\centering
\resizebox{\textwidth}{!}{
\begin{tikzpicture}[node distance=5mm]
  \node[box, fill=inkblue!12, minimum width=3.4cm] (root) {\textbf{59 automations}};
  \node[box, below left=11mm and 33mm of root, text width=3.0cm] (a) {\textbf{Lead gen \&\\outreach}\\\scriptsize 11};
  \node[box, right=6mm of a, text width=3.0cm] (b) {\textbf{Content\\generation}\\\scriptsize 14};
  \node[box, right=6mm of b, text width=3.0cm] (c) {\textbf{Video\\production}\\\scriptsize 7};
  \node[box, right=6mm of c, text width=3.0cm] (d) {\textbf{SEO}\\\scriptsize 5};
  \node[box, below=14mm of a, text width=3.0cm] (e) {\textbf{Assistants \&\\chatbots}\\\scriptsize 9};
  \node[box, right=6mm of e, text width=3.0cm] (f) {\textbf{Business ops}\\\scriptsize 6};
  \node[box, right=6mm of f, text width=3.0cm] (g) {\textbf{Utilities}\\\scriptsize 7};
  \node[extbox, right=6mm of g, text width=3.2cm] (h) {\scriptsize the dominant shape:\\\scriptsize trigger $\rightarrow$ fetch $\rightarrow$ LLM\\\scriptsize $\rightarrow$ deliver $\rightarrow$ log to Sheets};
  \foreach \x in {a,b,c,d} { \draw[flow] (root) -- (\x.north); }
  \foreach \x in {e,f,g} { \draw[flow] (root) to[out=270,in=90] (\x.north); }
\end{tikzpicture}}
\caption{Counts are of folders and sum to 59.}
\end{figure}

The extremes are worth knowing:

\begin{flushleft}
\begin{itemize}
  \item \textbf{Smallest automation:} 3 nodes, \file{youtube-transcript-fetcher}
    (a webhook, an API call, a reply).
  \item \textbf{Largest single file:} 190 nodes, \file{social-media-content-generator}.
  \item \textbf{Largest system:} 10 files and 275 nodes,
    \file{cold-email-outreach-sequence-system}.
\end{itemize}
\end{flushleft}

\section{Three automations worth understanding}

\subsection{The one that bills people}

\file{crm-lead-to-stripe-invoice} is 6 nodes and the best single example in the
repository, because it is small enough to hold in your head and its flaws cost real
money.

\begin{figure}[H]
\centering
\resizebox{\textwidth}{!}{
\begin{tikzpicture}[node distance=6mm]
  \node[extbox, text width=2.2cm] (cu) {ClickUp\\\scriptsize deal won};
  \node[box, right=8mm of cu, fill=warnred!12, draw=warnred, text width=1.8cm] (wh) {Webhook};
  \node[box, right=8mm of wh, text width=2.0cm] (get) {ClickUp\\\scriptsize get task};
  \node[box, right=8mm of get, text width=2.1cm] (st) {Stripe\\\scriptsize customer};
  \node[box, right=8mm of st, text width=2.0cm] (ii) {HTTP\\\scriptsize inv. item};
  \node[box, right=8mm of ii, text width=2.0cm] (ci) {HTTP\\\scriptsize invoice};
  \node[box, right=8mm of ci, text width=2.0cm] (fi) {HTTP\\\scriptsize finalize};
  \draw[flow] (cu) -- (wh); \draw[flow] (wh) -- (get); \draw[flow] (get) -- (st);
  \draw[flow] (st) -- (ii); \draw[flow] (ii) -- (ci); \draw[flow] (ci) -- (fi);
  \node[lbl, below=8mm of st, text width=8.6cm, draw=warnred, fill=warnred!8]
    {\textbf{verified in the node parameters:} the email is read as \code{custom\_fields[19]},
     the amount as \code{custom\_fields[5] * 100}, and the invoice line is the fixed
     string "1st Deposit". No error branch, no retry, no duplicate check anywhere in the chain.};
  \draw[dflow, draw=warnred] (get.south) -- ++(0,-4mm);
\end{tikzpicture}}
\caption{Six nodes, one straight line, no error path.}
\end{figure}

Note the shape: a native Stripe node creates the customer, then three hand-written HTTP
calls do the invoicing, because n8n's Stripe node does not implement Stripe's three-step
invoice dance. Native node where it fits, raw API where it does not, same stored
credential reused for both. That move is the defining habit of the whole collection.

\begin{honest}{Two real defects, both verified against the file}
\textbf{Fields are read by array position.} \code{custom\_fields[19]} is the email
\emph{today}. Reorder a column in ClickUp and it silently becomes something else. The
workflow does not fail; it invoices the wrong data to the wrong address, and the first
person to notice is the customer.

\textbf{No idempotency, on a node that charges people.} No check for an existing customer,
no check for an existing invoice. Webhook senders retry as a matter of routine. Two
deliveries produce two invoices. This is the highest-severity issue in the repository and
the only one that costs money on its first occurrence.
\end{honest}

\subsection{The one with the best architecture}

\file{ultimate-personal-assistant-system}: five files, 50 nodes. You text a Telegram bot,
by typing or by voice. An orchestrating agent works out what you want and delegates to one
of four specialist agents.

\begin{figure}[H]
\centering
\resizebox{0.93\textwidth}{!}{
\begin{tikzpicture}[node distance=6mm]
  \node[box, fill=warnred!12, draw=warnred, text width=2.0cm] (tg) {Telegram\\Trigger};
  \node[dec, right=8mm of tg] (sw) {text or\\voice?};
  \node[box, above right=3mm and 8mm of sw, text width=1.8cm] (set) {Set 'Text'};
  \node[box, below right=3mm and 8mm of sw, text width=1.8cm] (dl) {Download};
  \node[box, right=6mm of dl, text width=1.9cm] (tr) {Transcribe};
  \node[box, right=20mm of set, fill=accent!18, minimum height=14mm, text width=2.3cm] (ag) {\textbf{Ultimate}\\\textbf{Assistant}};
  \draw[flow] (tg) -- (sw);
  \draw[flow] (sw) -- node[lbl,above,pos=.4]{voice} (dl);
  \draw[flow] (sw) -- node[lbl,above,pos=.4]{text} (set);
  \draw[flow] (dl) -- (tr);
  \draw[flow] (set) -| (ag);
  \draw[flow] (tr) -| (ag);
  \node[box, below=12mm of ag, text width=2.0cm] (mem) {Memory};
  \node[box, left=4mm of mem, text width=2.0cm] (llm) {Chat Model};
  \draw[dflow] (llm) -- (ag); \draw[dflow] (mem) -- (ag);
  \node[box, above right=9mm and 12mm of ag, text width=2.2cm] (t1) {Email Agent};
  \node[box, below=2mm of t1, text width=2.2cm] (t2) {Calendar Agent};
  \node[box, below=2mm of t2, text width=2.2cm] (t3) {Contact Agent};
  \node[box, below=2mm of t3, text width=2.2cm] (t4) {Content Creator};
  \node[extbox, below=2mm of t4, text width=2.2cm] (t5) {Tavily / Calculator};
  \foreach \t in {t1,t2,t3,t4,t5} { \draw[flow, draw=accent] (ag) -- (\t); }
  \node[lbl, right=2mm of t2, text width=3.0cm, align=left] {each of these four is a\\\textbf{separate workflow file},\\called as a tool};
\end{tikzpicture}}
\caption{Solid blue arrows out of the agent are tool calls the model chooses to make.
Dashed arrows are equipment attached to it.}
\end{figure}

\begin{keyidea}{Sub-workflow as tool is the automation equivalent of a function call}
Each specialist is its own file with its own entry point, which buys what functions buy in
code: each can be tested alone, a change to contacts cannot break email, and each agent
carries a small focused prompt instead of one enormous one. The voice branch converges
back before the agent, so the agent never knows which input path it came from. That is the
right place for that seam.

The contrast is instructive: \file{social-media-content-generator} solves a comparable
problem on a single 190-node canvas.
\end{keyidea}

\subsection{The one with the most interesting internals}

\file{rag-ingestion-pipeline}, 47 nodes. It watches a Google Drive folder, and when a
document appears it works out the type, extracts the text, chops it up, converts the
chunks to embeddings and stores them in a Supabase vector table. A separate chat agent
answers questions from that table.

\begin{jargon}{RAG (Retrieval-Augmented Generation)}
A way to make an AI model answer from \emph{your} documents rather than its general
training. Chop the documents into chunks, convert each into a list of numbers capturing
its meaning (an \emph{embedding}), and store them. At question time, embed the question,
find the nearest chunks, and paste them into the prompt. The model then answers from real
text you supplied.
\end{jargon}

Two details show genuine understanding. Spreadsheets are summarized into prose
\emph{before} embedding, because embedding a spreadsheet row by row produces chunks like
\code{"Q3 | 41 | 12"} that carry no retrievable meaning. And on the update path, old rows
are deleted before new ones are inserted, because otherwise an edited document leaves its
old chunks behind and the agent answers from a version that no longer exists.

\begin{honest}{Half of it is a copy of the other half, and that half is switched off}
The file contains two complete, near-identical ingestion pipelines: one driven by a
\code{File Created} trigger, one by a \code{File Updated} trigger that is
\textbf{disabled}. Same nodes, different names (\code{Switch} and \code{Switch2},
\code{Aggregate} and \code{Aggregate1}).

Every future change to the extraction logic must be made twice, in two places, with
confusingly similar names, and a mistake in the dead copy stays invisible until someone
re-enables it. This is copy-paste architecture, and it is the characteristic debt of
visual tools: the canvas makes duplicating a branch trivially easy and offers nothing that
would encourage factoring it out.
\end{honest}

\section{The habits worth defending}

\subsection{Google Sheets as the database}

\begin{keyidea}{Defensible, with a known ceiling}
The instinct is to call it amateurish. A spreadsheet buys four things a real database does
not: the client can see it, the client can edit it (human review, for free), it needs no
provisioning, and it was already the incumbent tool. For a small agency automating its own
chores, that is an engineering decision, not a shortcut.

It breaks on \textbf{concurrency} (no transactions, so two runs can both read "not yet
emailed" and both email), on \textbf{uniqueness} (no constraint, which is precisely why the
cold email system can push a lead twice), and on \textbf{scale} (one SEO workflow reads its
entire post history on every run). Right call at this size; the ceiling is real and stated.
\end{keyidea}

\subsection{Trusting the AI's output shape}

The most common bug class here. A prompt says "return only JSON", and nothing downstream
checks that it did. n8n has a structured output parser node, which forces the response
into a declared schema and fails loudly when it does not fit. It appears \textbf{3 times
in 1{,}645 nodes}, in just 2 of the 80 files. In the other 78, output shape is a matter
of hoping.

\begin{honest}{The verified worst case, and the lesson in it}
In \file{competitor-content-gap-analyzer}, the \code{Clean Content} node's stored
JavaScript has had its newlines collapsed into literal letter \code{n} characters, so the
code reads \code{\{nnnn\$input.item.json.content = ...\}}. That throws a
\code{ReferenceError} the instant it runs. Its regular expressions are also passed as
quoted strings rather than regex literals, so they would match nothing even if the syntax
were fixed. The node has never cleaned anything.

The part that matters is what happens next: the node is set to
\code{continueOnFail: true}, so the error is swallowed, the item passes through
untouched, and \textbf{the workflow reports success}. Every run has produced a content gap
analysis based on no content.

The lesson is not that someone mangled a code node. It is that an error-tolerance flag
turned a loud crash into a silent wrong answer. That flag is set on 27 nodes here.
\end{honest}

\section{The honest assessment}
\label{sec:weak}

\subsection{Genuinely good}

\begin{itemize}
  \item \textbf{Documentation discipline is exceptional}: 59 READMEs on nine consistent
    headings, 59 credential manifests, 9{,}304 lines of prose, plus 186 sticky notes
    inside the canvases.
  \item \textbf{The self-criticism is real and specific.} The READMEs say "the
    content-cleaning code is broken" and "custom fields are read by array index".
    Spot-checks against the raw JSON confirmed these claims exactly, including the subtle
    ones. This is rarer than it should be.
  \item \textbf{No credentials leaked.} See below.
  \item \textbf{The hard parts were not dodged}: 128 hand-written API calls, async
    polling, a considered RAG chunking strategy, a properly factored sub-workflow
    architecture where it is used.
\end{itemize}

\subsection{Weak, ordered by what it would cost}

\begin{enumerate}
  \item \textbf{No idempotency on the invoicing workflow.} A retried webhook double-bills
    a real customer.
  \item \textbf{Silent failures via \code{continueOnFail}} (27 nodes). Proven, not
    theoretical.
  \item \textbf{Copy-paste architecture}: the duplicated RAG branch; six duplicated
    copywriting blocks in the cold email system; nine orphaned call nodes left from an
    old version.
  \item \textbf{Orphaned nodes in 12 of the 80 files}, 55-plus nodes with no path from any
    trigger. On a canvas they look exactly like working steps.
  \item \textbf{Unbounded polling loops.} No attempt cap anywhere in the video family.
  \item \textbf{Hardcoded values where expressions belong}: a fixed \code{\$1,850} cost,
    "Hey Nick" in a client email, the literal caption \code{"hello"} on every post in one
    video workflow.
  \item \textbf{Documentation drift}: sticky notes labelling reminder stages "5-Minute
    Alert" when the node conditions say one hour.
\end{enumerate}

\subsection{Security}

\begin{keyidea}{No live credential is committed, and this was checked properly}
n8n stores logins in an encrypted credential store; a node holds only a pointer to one.
Verified across all 80 files: \textbf{373 credential references}, 27 distinct credential
types, every one containing exactly \code{id} and \code{name}. Not one carries a value.

That is n8n's design doing the work. What the author was responsible for is everything
n8n does not protect, and that was audited too: all provider-key formats scanned across
all 268 tracked files (three hits, all placeholders such as
\code{Bearer YOUR\_API\_KEY\_HERE}); every referenced host is a public vendor API, with no
instance hostname and no Supabase project URL; every email is a placeholder; the
\code{shared} block that normally carries the account owner's name and email is scrubbed
in all 30 files that have one.

The project README says live keys were found and stripped. This audit found no survivor.
\end{keyidea}

\begin{honest}{The scrub was manual, and that is the residual risk}
One field, the n8n account name, appears with \textbf{sixteen different placeholder
spellings} across the files (\code{YOUR\_EMAIL\_HERE}, \code{owner@example.com},
\code{Portfolio Owner}, \code{REDACTED\_USER}, and twelve more). The outcome is correct,
but sixteen spellings of one substitution is the signature of a hand pass, and a hand pass
is one that can miss a file. A scripted scrub with a single canonical placeholder would be
safer, because it could be re-run to \emph{prove} nothing was missed. Today the proof is
an audit, which is only a snapshot.
\end{honest}

\subsection{The maturity question}

\begin{honest}{The one thing to get straight before any interview}
\textbf{28 manual triggers} (the most common trigger type here). \textbf{9 of 80}
workflows marked active. \textbf{45} disabled nodes. \textbf{12} files with orphans.

A manual trigger is what n8n gives you while you are \emph{building}, so you can press
Execute and watch data flow. A deployed automation is normally started by a schedule, a
webhook or an event.

The defensible claim is: \emph{59 automations designed and built across seven problem
domains, every one documented including its faults, with a clear account of which are
production-hardened and which are at demonstration maturity.} That is a strong claim and
it survives someone opening the files.

The claim to avoid is "59 automations run my client's business". The repository does not
contain the evidence for it: execution logs are not here, and the active flags say what
they say.
\end{honest}

\subsection{What is inferred rather than confirmed}

The seven-family taxonomy is this document's own reading; the repository does not group
its folders, and while the folder counts are exact the boundaries are a judgment call. The
claim that the collection is "built and demonstrated rather than deployed" is a strong
inference from trigger types and active flags, but the alternative (that workflows were
routinely deactivated before export) cannot be ruled out from the files alone. In the cold
email system, the handoff from the research workflow to the orchestrator is inferred from
matching field names; there is no connection in the JSON, and the project's own README is
equally explicit about that. Whether the disabled trigger in the RAG pipeline was switched
off deliberately or abandoned mid-refactor cannot be determined from the export. And
whether the assistants behave \emph{well} is untestable from the files: the graphs are
sound, but the behaviour lives in prompts, and prompts cannot be evaluated by reading them.

\vfill
\begin{center}
\rule{0.4\textwidth}{0.4pt}\\[0.3em]
{\small All counts were produced by parsing the 80 exported workflow files directly. The
full treatment, including the data model of an export, walkthroughs of the largest
systems, and the complete audit, is in the companion Deep Dive.}
\end{center}

\end{document}
